\documentclass[12pt]{article}
\usepackage{amsmath, amssymb}
\usepackage{geometry}
\usepackage{enumitem}
\geometry{margin=1in}

\title{Algebra II Logs and Series Practice Test}
\author{}
\date{}

\begin{document}
\maketitle

\section*{Instructions}
Complete in 60 minutes. Show clear work. Use exact values unless a decimal approximation is explicitly requested. For logarithmic equations, state domain restrictions and check for extraneous solutions.

\section*{Reference Facts}
\begin{itemize}[leftmargin=1.5em]
  \item Log properties:
  \[
  \log_b(MN)=\log_bM+\log_bN,\qquad
  \log_b\!\left(\frac{M}{N}\right)=\log_bM-\log_bN,\qquad
  \log_b(M^r)=r\log_bM
  \]
  \[
  \log_b(b^x)=x,\qquad b^{\log_b x}=x,\qquad
  \log_bx=\frac{\log_a x}{\log_a b}
  \]
  \item Arithmetic and geometric sequences:
  \[
  a_n=a_1+(n-1)d,\qquad S_n=\frac{n(a_1+a_n)}{2}
  \]
  \[
  g_n=g_1r^{n-1},\qquad S_n=\frac{g_1(1-r^n)}{1-r},\qquad
  S_\infty=\frac{g_1}{1-r}\quad(|r|<1)
  \]
  \item Consecutive integer sum:
  \[
  1+2+3+\cdots+n=\frac{n(n+1)}{2}
  \]
\end{itemize}

\section*{Problems}
\begin{enumerate}[leftmargin=*, itemsep=1.1em, topsep=0.8em]
  \item Determine whether the relation is a function. Then determine whether its inverse relation is a function. Explain briefly.
  \[
  (-4,9),\ (-2,3),\ (0,-1),\ (2,3),\ (4,9)
  \]

  \item Let
  \[
  f(x)=\frac{x+6}{x-2}.
  \]
  Find the domain, range, $x$-intercept, $y$-intercept, and both asymptotes.

  \item Determine whether each function is even, odd, or neither. Justify algebraically.
  \[
  f(x)=x^5-4x^3+x,\qquad
  g(x)=x^4-3x^2+8,\qquad
  h(x)=x^3+x^2-2x.
  \]

  \item Expand completely:
  \[
  \log_b\!\left(\frac{16x^3\sqrt{y}}{b^4\sqrt[3]{z^5}}\right).
  \]

  \item Condense into a single logarithm:
  \[
  2\log_5(x-1)-\frac12\log_5(y+3)+3\log_5(25)-\log_5(z^2).
  \]

  \item Solve and check for extraneous solutions:
  \[
  \log_3(x-2)+\log_3(x-6)=2.
  \]

  \item Solve and check for extraneous solutions:
  \[
  \log_2(x+4)-\log_2(x-1)=\log_2(x-2).
  \]

  \item Evaluate exactly.
  \[
  \sum_{k=1}^{30}(4k-7),\qquad
  \sum_{j=8}^{42}j
  \]

  \item For the arithmetic sequence with $a_5=19$ and $a_{18}=71$, find $a_1$, $d$, and $S_{25}$.

  \item For the geometric sequence with $g_3=48$ and $g_7=3$, find all possible values of $r$ and $g_1$. Then find $g_{10}$ for each possible sequence.

  \item Evaluate exactly.
  \[
  \sum_{n=0}^{7}6\left(\frac23\right)^n,\qquad
  \sum_{k=1}^{n}(5k-2)
  \]

  \item Determine whether each infinite geometric series converges. If it converges, find its exact sum.
  \[
  \sum_{n=1}^{\infty}18\left(-\frac13\right)^{n-1},\qquad
  \sum_{n=0}^{\infty}5\left(\frac43\right)^n
  \]
\end{enumerate}

\end{document}
